\section{Partition basics}

\subsection{Definitions}

\begin{definition}
\label{def.pars.parts}
\lean{Nat.Partition, Nat.Partition.size_eq, Nat.Partition.parts_pos', Nat.Partition.ofList', Nat.Partition.eq_iff_parts_eq}
\leantarget
\leanok
\textbf{(a)} An \emph{(integer) partition} means a
(finite) weakly decreasing tuple of positive integers -- i.e., a finite tuple
$\left(  \lambda_{1},\lambda_{2},\ldots,\lambda_{m}\right)  $ of positive
integers such that $\lambda_{1}\geq\lambda_{2}\geq\cdots\geq\lambda_{m}$.

Thus, partitions are the same as weakly decreasing compositions. Hence, the
notions of \emph{size} and \emph{length} of a partition are automatically
defined, since we have defined them for compositions (in Definition
\ref{def.fps.comps}). \medskip

\textbf{(b)} The \emph{parts} of a partition $\left(  \lambda_{1},\lambda
_{2},\ldots,\lambda_{m}\right)  $ are simply its entries $\lambda_{1}%
,\lambda_{2},\ldots,\lambda_{m}$. \medskip

\textbf{(c)} Let $n\in\mathbb{Z}$. A \emph{partition of }$n$ means a partition
whose size is $n$. \medskip

\textbf{(d)} Let $n\in\mathbb{Z}$ and $k\in\mathbb{N}$. A \emph{partition of
}$n$\emph{ into }$k$\emph{ parts} is a partition whose size is $n$ and whose
length is $k$.
\end{definition}

\begin{lemma}
\label{lem.pars.parts-le}
\lean{Nat.Partition.parts_le}
\leanhelper
\leanok
A partition of $n$ has all parts $\leq n$.
\end{lemma}

\begin{proof}
\leanok
Since parts are positive and sum to $n$, each individual part is bounded by $n$.
\end{proof}

\begin{lemma}
\label{lem.pars.parts-zero}
\lean{Nat.Partition.parts_eq_zero_of_partition_zero}
\leanhelper
\leanok
The partition of $0$ has no parts (i.e., it is the empty tuple).
\end{lemma}

\begin{proof}
\leanok
By definition.
\end{proof}

\begin{lemma}
\label{lem.pars.parts-card-zero}
\lean{Nat.Partition.parts_card_zero}
\leanhelper
\leanok
The length of the partition of $0$ is $0$.
\end{lemma}

\begin{proof}
\leanok
Follows from the fact that the partition of $0$ has no parts.
\end{proof}

\begin{lemma}
\label{lem.pars.card-le-sum-of-pos}
\lean{Nat.Partition.card_le_sum_of_pos}
\leanhelper
\leanok
The cardinality of a multiset of positive naturals is at most its sum.
\end{lemma}

\begin{proof}
\leanok
By induction on the multiset.
\end{proof}

\begin{lemma}
\label{lem.pars.numParts-le}
\lean{Nat.Partition.numParts_le}
\leanhelper
\leanok
A partition of $n$ has at most $n$ parts.
\end{lemma}

\begin{proof}
\leanok
Since each part is $\geq 1$, the number of parts is at most the sum of parts, which is $n$.
\end{proof}

\begin{lemma}
\label{lem.pars.numParts-pos}
\lean{Nat.Partition.numParts_pos}
\leanhelper
\leanok
For $n>0$, every partition of $n$ has at least one part.
\end{lemma}

\begin{proof}
\leanok
If a partition of a positive $n$ had no parts, its sum would be~$0$, a contradiction.
\end{proof}

\begin{lemma}
\label{lem.pars.largestPart-pos}
\lean{Nat.Partition.largestPart_pos}
\leanhelper
\leanok
For $n>0$, every partition of $n$ has a positive largest part.
\end{lemma}

\begin{proof}
\leanok
Since the partition has at least one part and all parts are positive, the maximum is positive.
\end{proof}

\begin{lemma}
\label{lem.pars.partitionCount-pos}
\lean{Nat.Partition.partitionCount_pos}
\leanhelper
\leanok
For every $n\in\mathbb{N}$, we have $p(n) > 0$.
\end{lemma}

\begin{proof}
\leanok
The empty partition (for $n=0$) or the indiscrete partition $(n)$ (for $n>0$)
provides at least one partition.
\end{proof}

\begin{definition}
\label{def.pars.pn-pkn}
\lean{Nat.Partition.partitionCount, Nat.Partition.partsCount, Nat.Partition.partsLeqCount, Nat.Partition.partsInCount}
\leantarget
\leanok
\textbf{(a)} Let $n\in\mathbb{Z}$ and $k\in\mathbb{N}$.
Then, we set%
\[
p_{k}\left(  n\right)  :=\left(  \text{\# of partitions of }n\text{ into
}k\text{ parts}\right)  .
\]

\textbf{(b)} Let $n\in\mathbb{Z}$. Then, we set%
\[
p\left(  n\right)  :=\left(  \text{\# of partitions of }n\right)  .
\]
This is called the $n$\emph{-th partition number}.
\end{definition}

\begin{definition}
\label{def.pars.iverson}
\lean{IversonBracket.iverson, IversonBracket.iverson_true, IversonBracket.iverson_false, IversonBracket.kronecker_eq_iverson}
\leantarget
\leanok
We will use the \emph{Iverson bracket notation}: If
$\mathcal{A}$ is a logical statement, then $\ive{\mathcal{A}}$
means the \emph{truth value} of $\mathcal{A}$; this is the integer $%
\begin{cases}
1, & \text{if }\mathcal{A}\text{ is true};\\
0, & \text{if }\mathcal{A}\text{ is false}.
\end{cases}
$

Note that the Kronecker delta notation is a particular case of the Iverson
bracket: We have
\[
\delta_{i,j}=\ive{i=j}\ \ \ \ \ \ \ \ \ \ \text{for any objects
}i\text{ and }j.
\]
\end{definition}

\begin{lemma}
\label{lem.pars.iverson-eq-one-iff}
\lean{IversonBracket.iverson_eq_one_iff}
\leanhelper
\leanok
The Iverson bracket equals $1$ iff the proposition is true.
\end{lemma}

\begin{proof}
\leanok
By case analysis on whether the proposition is true or false.
\end{proof}

\begin{lemma}
\label{lem.pars.iverson-eq-zero-iff}
\lean{IversonBracket.iverson_eq_zero_iff}
\leanhelper
\leanok
The Iverson bracket equals $0$ iff the proposition is false (assuming $1 \neq 0$).
\end{lemma}

\begin{proof}
\leanok
By case analysis on whether the proposition is true or false.
\end{proof}

\begin{lemma}
\label{lem.pars.sum-iverson-eq-card}
\lean{IversonBracket.sum_iverson_eq_card}
\leanhelper
\leanok
The sum of Iverson brackets equals the cardinality of the filtered set:
$\sum_{x \in s} \ive{p(x)} = |s \cap \{x : p(x)\}|$.
\end{lemma}

\begin{proof}
\leanok
Follows from the definition of the Iverson bracket and the identity
$\sum_{x \in s} \ive{p(x)} = |\{x \in s : p(x)\}|$.
\end{proof}

\begin{lemma}
\label{lem.pars.iverson-and}
\lean{IversonBracket.iverson_and}
\leanhelper
\leanok
The Iverson bracket is multiplicative for conjunctions:
$\ive{P \wedge Q} = \ive{P} \cdot \ive{Q}$.
\end{lemma}

\begin{proof}
\leanok
By case split on both propositions.
\end{proof}

\begin{definition}
\label{def.pars.floor-ceil}
\lean{FloorCeiling.floor_def, FloorCeiling.ceil_def, FloorCeiling.floor_of_int, FloorCeiling.ceil_of_int, FloorCeiling.nat_div_eq_floor}
\leantarget
\leanok
Let $a$ be a real number.

Then, $\left\lfloor a\right\rfloor $ (called the \emph{floor} of $a$) means
the largest integer that is $\leq a$.

Likewise, $\left\lceil a\right\rceil $ (called the \emph{ceiling} of $a$)
means the smallest integer that is $\geq a$.
\end{definition}

\subsection{Simple properties of partition numbers}

\begin{convention}
\label{conv.pars.largest-part-0}
\lean{Nat.Partition.largestPart}
\leanok
We agree to say that the largest part of the
empty partition $\left(  {}\right)  $ is $0$ (even though this partition has
no parts).
\end{convention}

\begin{lemma}
\label{lem.pars.largestPart-zero}
\lean{Nat.Partition.largestPart_zero}
\leanhelper
\leanok
The largest part of the partition of $0$ is $0$.
\end{lemma}

\begin{proof}
\leanok
The partition of $0$ has no parts, so the maximum (with default $0$) is $0$.
\end{proof}

\begin{lemma}
\label{lem.pars.largestPart-le}
\lean{Nat.Partition.largestPart_le}
\leanhelper
\leanok
For a partition of $n$, the largest part is at most $n$.
\end{lemma}

\begin{proof}
\leanok
Each part is bounded by $n$ (by Lemma~\ref{lem.pars.parts-le}).
\end{proof}

\begin{lemma}
\label{lem.pars.largestPart-le-iff}
\lean{Nat.Partition.largestPart_le_iff_all_parts_le}
\leanhelper
\leanok
A partition has largest part $\leq m$ if and only if all its parts are $\leq m$.
\end{lemma}

\begin{proof}
\leanok
The largest part is the maximum of all parts, so it is $\leq m$ iff every part is $\leq m$.
\end{proof}

\begin{lemma}
\label{lem.pars.filter-largestPart-le-eq-restricted}
\lean{Nat.Partition.filter_largestPart_le_eq_restricted}
\leanhelper
\leanok
The set of partitions of $n$ whose largest part is $\le m$ equals
the set of partitions of $n$ with all parts $\le m$.
\end{lemma}

\begin{proof}
Follows directly from Lemma~\ref{lem.pars.largestPart-le-iff}.
\end{proof}

\begin{proposition}
\label{prop.pars.basics}
\lean{Nat.Partition.partsCount_of_gt, Nat.Partition.partsCount_zero, Nat.Partition.partsCount_one, Nat.Partition.partsCount_recurrence, Nat.Partition.partsCount_two, Nat.Partition.partitionCount_sum}
\leantarget
\leanok
Let $n\in\mathbb{Z}$ and $k\in\mathbb{N}$. \medskip

\textbf{(a)} We have $p_{k}\left(  n\right)  =0$ whenever $n<0$ and
$k\in\mathbb{N}$. \medskip

\textbf{(b)} We have $p_{k}\left(  n\right)  =0$ whenever $k>n$. \medskip

\textbf{(c)} We have $p_{0}\left(  n\right)  =\ive{n=0}$. \medskip

\textbf{(d)} We have $p_{1}\left(  n\right)  =\ive{n>0}$. \medskip

\textbf{(e)} We have $p_{k}\left(  n\right)  =p_{k}\left(  n-k\right)
+p_{k-1}\left(  n-1\right)  $ whenever $k>0$. \medskip

\textbf{(f)} We have $p_{2}\left(  n\right)  =\left\lfloor n/2\right\rfloor $
whenever $n\in\mathbb{N}$. \medskip

\textbf{(g)} We have $p\left(  n\right)  =p_{0}\left(  n\right)  +p_{1}\left(
n\right)  +\cdots+p_{n}\left(  n\right)  $ whenever $n\in\mathbb{N}$. \medskip

\textbf{(h)} We have $p\left(  n\right)  =0$ whenever $n<0$.
\end{proposition}

\begin{proof}
\textbf{(a)} The size
of a partition is always nonnegative (being a sum of positive integers). Thus,
a negative number $n$ has no partitions whatsoever. Thus, $p_{k}\left(
n\right)  =0$ whenever $n<0$ and $k\in\mathbb{N}$.
(In the formalization, $n$ is a natural number, so this case is vacuous.) \medskip

\textbf{(b)} If $\left(  \lambda_{1},\lambda_{2},\ldots,\lambda_{k}\right)  $
is a partition, then $\lambda_{i}\geq1$ for each $i\in\left\{  1,2,\ldots
,k\right\}  $ (because a partition is a tuple of positive integers, i.e., of
integers $\geq1$). Hence, if $\left(  \lambda_{1},\lambda_{2},\ldots
,\lambda_{k}\right)  $ is a partition of $n$ into $k$ parts, then%
\begin{align*}
n  &  =\lambda_{1}+\lambda_{2}+\cdots+\lambda_{k}\geq\underbrace{1+1+\cdots+1}_{k\text{ times}}=k.
\end{align*}
Thus, a partition of $n$ into $k$ parts cannot satisfy $k>n$. Thus, no such
partitions exist if $k>n$. In other words, $p_{k}\left(  n\right)  =0$ if
$k>n$. \medskip

\textbf{(c)} The integer $0$ has a unique partition into $0$ parts, namely the
empty tuple $\left(  {}\right)  $. A nonzero integer $n$ cannot have any
partitions into $0$ parts, since the empty tuple has size $0\neq n$. Thus,
$p_{0}\left(  n\right)  $ equals $1$ for $n=0$ and equals $0$ for $n\neq0$. In
other words, $p_{0}\left(  n\right)  =\ive{n=0}$. \medskip

\textbf{(d)} Any positive integer $n$ has a unique partition into $1$ part --
namely, the $1$-tuple $\left(  n\right)  $. On the other hand, if $n$ is not
positive, then this $1$-tuple is not a partition, so in this case $n$ has no
partition into $1$ part. Thus, $p_{1}\left(  n\right)  $ equals $1$ if $n$ is
positive and equals $0$ otherwise. In other words, $p_{1}\left(  n\right)
=\ive{n>0}$. \medskip

\textbf{(e)} Assume that $k>0$. We must prove that $p_{k}\left(  n\right)
=p_{k}\left(  n-k\right)  +p_{k-1}\left(  n-1\right)  $.

We consider all partitions of $n$ into $k$ parts. We classify these partitions
into two types:

\begin{itemize}
\item \emph{Type 1} consists of all partitions that have $1$ as a part.

\item \emph{Type 2} consists of all partitions that don't.
\end{itemize}

Any type-1 partition has $1$ as a part, therefore as its last part (because it
is weakly decreasing). Hence, any type-1 partition has the form $\left(
\lambda_{1},\lambda_{2},\ldots,\lambda_{k-1},1\right)  $. If $\left(
\lambda_{1},\lambda_{2},\ldots,\lambda_{k-1},1\right)  $ is a type-1 partition
(of $n$ into $k$ parts), then $\left(  \lambda_{1},\lambda_{2},\ldots
,\lambda_{k-1}\right)  $ is a partition of $n-1$ into $k-1$ parts. Thus, we
have a bijection
\[
\left\{  \text{type-1 partitions of }n\text{ into }k\text{ parts}\right\}
\to\left\{  \text{partitions of }n-1\text{ into }k-1\text{ parts}\right\}
\]
given by removing the trailing $1$.

A type-2 partition does not have $1$ as a
part; hence, all its parts are larger than $1$, and therefore we can subtract $1$ from each part and still
have a partition. If $\left(  \lambda
_{1},\lambda_{2},\ldots,\lambda_{k}\right)  $ is a type-2 partition of $n$
into $k$ parts, then $\left(  \lambda_{1}-1,\lambda_{2}-1,\ldots,\lambda_{k}-1\right)  $
is a partition of $n-k$ into $k$ parts. This gives a bijection
\[
\left\{  \text{type-2 partitions of }n\text{ into }k\text{ parts}\right\}
\to\left\{  \text{partitions of }n-k\text{ into }k\text{ parts}\right\}.
\]

Since any partition of $n$ into $k$ parts is either type-1 or type-2,
$p_{k}\left(  n\right)  =p_{k}\left(  n-k\right)  +p_{k-1}\left(  n-1\right)$. \medskip

\textbf{(f)} Let $n\in\mathbb{N}$. The partitions of $n$ into $2$ parts are%
\[
\left(  n-1,1\right)  ,\ \ \left(  n-2,2\right)  ,\ \ \ldots,\ \ \left(
\left\lceil n/2\right\rceil ,\left\lfloor n/2\right\rfloor \right)  .
\]
Thus there are $\left\lfloor n/2\right\rfloor $ of them. In other words,
$p_{2}\left(  n\right)  =\left\lfloor n/2\right\rfloor $. \medskip

\textbf{(g)} Let $n\in\mathbb{N}$. Any partition of $n$ must have $k$ parts
for some $k\in\mathbb{N}$. Thus,%
\begin{align*}
p\left(  n\right)   &  =\sum_{k\in\mathbb{N}}p_{k}\left(  n\right)
=\sum_{k=0}^{n}p_{k}\left(  n\right)
=p_{0}\left(  n\right)  +p_{1}\left(
n\right)  +\cdots+p_{n}\left(  n\right)  .
\end{align*}

\textbf{(h)} Same argument as for \textbf{(a)}.
\end{proof}

\subsubsection{Helpers for the recurrence (Proposition~\ref{prop.pars.basics}~(e))}

The recurrence $p_k(n) = p_k(n-k) + p_{k-1}(n-1)$ requires two bijections.
The following definitions and lemmas provide the infrastructure.

\begin{definition}
\label{def.pars.removeOne-addOne}
\lean{Nat.Partition.removeOne, Nat.Partition.addOne}
\leanhelper
\leanok
Given a partition of $n$ that contains $1$ as a part, removing
one copy of~$1$ yields a partition of $n-1$.
Conversely, adding one copy of~$1$ to a partition of $n$
yields a partition of $n+1$.
\end{definition}

\begin{definition}
\label{def.pars.subtractOneFromEach-addOneToEach}
\lean{Nat.Partition.subtractOneFromEach, Nat.Partition.addOneToEach}
\leanhelper
\leanok
Given a partition of $n$ into $k$ parts with all parts $>1$,
subtracting~$1$ from each part yields a
partition of $n-k$ into $k$ parts.
Conversely, adding~$1$ to each part of a partition of $n$ into $k$ parts produces a
partition of $n+k$ into $k$ parts.
\end{definition}

\begin{lemma}
\label{lem.pars.partsCount-split}
\lean{Nat.Partition.partsCount_split}
\leanhelper
\leanok
The set of partitions of $n$ into $k$ parts splits into those containing~$1$
and those not containing~$1$:
\[
p_k(n) = |\{\text{partitions with 1}\}| + |\{\text{partitions without 1}\}|.
\]
\end{lemma}

\begin{proof}
By the disjoint union of the two filtered sets.
\end{proof}

\begin{lemma}
\label{lem.pars.partsWithOne-card-eq}
\lean{Nat.Partition.partsWithOne_card_eq}
\leanhelper
\leanok
For $k>0$ and $n>0$:
\[
|\{\text{partitions of $n$ into $k$ parts containing 1}\}| = p_{k-1}(n-1).
\]
The bijection removes/adds a trailing~$1$.
\end{lemma}

\begin{proof}
The forward map removes one copy of~$1$; the backward map adds one.
Injectivity and surjectivity follow from the fact that removing and adding a trailing~$1$
are mutually inverse operations.
\end{proof}

\begin{lemma}
\label{lem.pars.partsWithoutOne-card-eq}
\lean{Nat.Partition.partsWithoutOne_card_eq}
\leanhelper
\leanok
For $k>0$:
\[
|\{\text{partitions of $n$ into $k$ parts not containing 1}\}| = p_k(n-k).
\]
The bijection subtracts/adds~$1$ from/to each part.
\end{lemma}

\begin{proof}
The forward map subtracts~$1$ from each part (valid since all parts $>1$).
The backward map adds~$1$ to each part.
Injectivity uses the fact that subtracting~$1$ from each entry is injective on multisets
whose all elements are $>1$.
\end{proof}

\subsubsection{Helpers for $p_2(n) = \lfloor n/2 \rfloor$ (Proposition~\ref{prop.pars.basics}~(f))}

\begin{lemma}
\label{lem.pars.partition-two-parts-form}
\lean{Nat.Partition.partition_two_parts_form}
\leanhelper
\leanok
Every partition of $n$ into $2$ parts has the form $\{n-b,\,b\}$ for some
$1\le b \le \lfloor n/2\rfloor$.
\end{lemma}

\begin{proof}
If the partition has parts $\{a, b\}$ with $b\le a$, then $a+b = n$ and $b\ge 1$
(since parts are positive), so $a = n-b$ and $2b\le n$.
\end{proof}

\begin{lemma}
\label{lem.pars.partsCount-eq-zero-iff}
\lean{Nat.Partition.partsCount_eq_zero_iff}
\leanhelper
\leanok
We have $p_k(n) = 0$ if and only if $k > n$ or ($k = 0$ and $n > 0$).
\end{lemma}

\begin{proof}
\leanok
The forward direction follows by constructing an explicit partition with $k$ parts when $0 < k \leq n$
(namely $(1, 1, \ldots, 1, n - k + 1)$). The backward direction uses parts (b) and (c) of
Proposition~\ref{prop.pars.basics}.
\end{proof}

\subsection{The generating function}

\begin{theorem}
\label{thm.pars.main-gf}
\lean{Nat.Partition.partitionCount_genFun}
\leantarget
\leanok
In the FPS ring $\mathbb{Z}\left[  \left[  x\right]
\right]  $, we have%
\[
\sum_{n\in\mathbb{N}}p\left(  n\right)  x^{n}=\prod_{k=1}^{\infty}\dfrac
{1}{1-x^{k}}.
\]

(The product on the right hand side is well-defined, since multiplying a FPS
by $\dfrac{1}{1-x^{k}}$ does not affect its first $k$ coefficients.)
\end{theorem}

\begin{proof}
We have%
\begin{align*}
&  \prod_{k=1}^{\infty}\dfrac{1}{1-x^{k}}
=\prod_{k=1}^{\infty}\sum_{u\in\mathbb{N}}x^{ku}
=\sum_{\substack{\left(  u_{1},u_{2},u_{3},\ldots\right)  \in
\mathbb{N}^{\infty}\text{ is}\\\text{essentially finite}}}x^{1u_{1}%
+2u_{2}+3u_{3}+\cdots}
=\sum_{n\in\mathbb{N}}\left\vert Q_{n}\right\vert x^{n},
\end{align*}
where
\[
Q_{n}=\left\{  \left(  u_{1},u_{2},u_{3},\ldots\right)  \in\mathbb{N}^{\infty
}\text{ essentially finite }\mid\ 1u_{1}+2u_{2}+3u_{3}+\cdots=n\right\}  .
\]

It suffices to show that $\left\vert Q_{n}\right\vert =p\left(  n\right)  $ for each $n\in\mathbb{N}$.

We construct a bijection $\pi:Q_{n}\rightarrow\left\{  \text{partitions of }n\right\}  $:
for any $\left(  u_{1},u_{2},u_{3},\ldots\right)  \in Q_{n}$, define
\begin{align*}
\pi\left(  u_{1},u_{2},u_{3},\ldots\right)  :=  &  \left(  \text{the partition
that contains each }i\text{ exactly }u_{i}\text{ times}\right) \\
=  &  \left(  \ldots,\underbrace{3,3,\ldots,3}_{u_{3}\text{ times}%
},\underbrace{2,2,\ldots,2}_{u_{2}\text{ times}},\underbrace{1,1,\ldots
,1}_{u_{1}\text{ times}}\right)  .
\end{align*}
This is a partition of $n$. Its inverse map $\rho$ sends each partition $\lambda$ of $n$ to
$\left(  \text{\# of }1\text{'s in }\lambda,\text{ \# of }2\text{'s in
}\lambda, \ldots\right)  \in Q_{n}$.
\end{proof}

\begin{theorem}
\label{thm.pars.main-gf-parts-n}
\lean{Nat.Partition.partitionCount_genFun_partsLeq, Nat.Partition.partitionCount_genFun_partsLeq_finprod}
\leantarget
\leanok
Let $m\in\mathbb{N}$. For each $n\in
\mathbb{N}$, let $p_{\operatorname{parts}\leq m}\left(  n\right)  $ be the \#
of partitions $\lambda$ of $n$ such that all parts of $\lambda$ are $\leq m$.
Then,
\[
\sum_{n\in\mathbb{N}}p_{\operatorname{parts}\leq m}\left(  n\right)
x^{n}=\prod_{k=1}^{m}\dfrac{1}{1-x^{k}}.
\]
\end{theorem}

\begin{proof}
This proof is analogous to the proof of Theorem~\ref{thm.pars.main-gf},
using $m$-tuples instead of infinite sequences.
We have
\[
\prod_{k=1}^{m}\dfrac{1}{1-x^{k}}=\prod_{k=1}^{m}\sum_{u\in
\mathbb{N}}x^{ku}=\sum_{\left(  u_{1},u_{2},\ldots,u_{m}\right)  \in\mathbb{N}^{m}}%
x^{1u_{1}+2u_{2}+\cdots+mu_{m}}=\sum_{n\in\mathbb{N}}\left\vert Q_{n}%
\right\vert x^{n},
\]
where
$Q_{n}=\left\{  \left(  u_{1},u_{2},\ldots,u_{m}\right)  \in\mathbb{N}%
^{m} \mid\ 1u_{1}+2u_{2}+\cdots+mu_{m}=n\right\}$.
It suffices to show $\left\vert Q_{n}\right\vert =p_{\operatorname{parts}\leq m}\left(  n\right)$.
The bijection $\pi:Q_{n}\rightarrow\left\{  \text{partitions }\lambda\text{ of }n\text{ with
all parts }\leq m\right\}  $ is analogous to that in Theorem~\ref{thm.pars.main-gf}.
\end{proof}

\begin{theorem}
\label{thm.pars.main-gf-parts-I}
\lean{Nat.Partition.partitionCount_genFun_partsIn}
\leantarget
\leanok
Let $I$ be a subset of $\left\{
1,2,3,\ldots\right\}  $. For each $n\in\mathbb{N}$, let $p_{I}\left(
n\right)  $ be the \# of partitions $\lambda$ of $n$ such that all parts of
$\lambda$ belong to $I$. Then,
\[
\sum_{n\in\mathbb{N}}p_{I}\left(  n\right)  x^{n}=\prod_{k\in I}\dfrac
{1}{1-x^{k}}.
\]
\end{theorem}

\begin{proof}
Analogous to the proof of Theorem~\ref{thm.pars.main-gf-parts-n}, with
$\prod_{k=1}^{m}$ replaced by $\prod_{k\in I}$ and
$m$-tuples replaced by essentially finite families $(u_i)_{i\in I}\in\mathbb{N}^{I}$.
\end{proof}

\subsection{Odd parts and distinct parts}

\begin{definition}
\label{def.pars.odd-dist-parts}
\lean{Nat.Partition.IsOddParts, Nat.Partition.IsDistinctParts, Nat.Partition.oddPartsCount, Nat.Partition.distinctPartsCount}
\leantarget
\leanok
Let $n\in\mathbb{Z}$. \medskip

\textbf{(a)} A \emph{partition of }$n$\emph{ into odd parts} means a partition
of $n$ whose all parts are odd. \medskip

\textbf{(b)} A \emph{partition of }$n$\emph{ into distinct parts} means a
partition of $n$ whose parts are distinct. \medskip

\textbf{(c)} Let%
\begin{align*}
p_{\operatorname{odd}}\left(  n\right)   &  :=\left(  \text{\# of partitions
of }n\text{ into odd parts}\right)  \ \ \ \ \ \ \ \ \ \ \text{and}\\
p_{\operatorname{dist}}\left(  n\right)   &  :=\left(  \text{\# of partitions
of }n\text{ into distinct parts}\right)  .
\end{align*}
\end{definition}

\begin{lemma}
\label{lem.pars.isOddParts-iff}
\lean{Nat.Partition.isOddParts_iff}
\leanhelper
\leanok
A partition is into odd parts iff all its parts are odd.
\end{lemma}

\begin{proof}
\leanok
By definition.
\end{proof}

\begin{lemma}
\label{lem.pars.isDistinctParts-iff}
\lean{Nat.Partition.isDistinctParts_iff}
\leanhelper
\leanok
A partition is into distinct parts iff its parts multiset has no duplicates.
\end{lemma}

\begin{proof}
\leanok
By definition.
\end{proof}

\begin{lemma}
\label{lem.pars.isDistinctParts-iff-count}
\lean{Nat.Partition.isDistinctParts_iff_count_le_one}
\leanhelper
\leanok
A partition is into distinct parts iff each part appears at most once.
\end{lemma}

\begin{proof}
\leanok
This is the characterization of ``no duplicates'' in terms of counts.
\end{proof}

\begin{theorem}[Euler's odd-distinct identity]
\label{thm.pars.odd-dist-equal}
\lean{Nat.Partition.odd_eq_distinct, Nat.Partition.oddPartsCount_eq_distinctPartsCount}
\leantarget
\leanok
We have
$p_{\operatorname{odd}}\left(  n\right)  =p_{\operatorname{dist}}\left(
n\right)  $ for each $n\in\mathbb{N}$.
\end{theorem}

\begin{proof}
\leanok
Let $n\in\mathbb{N}$. We construct a bijection
\[
A:\left\{  \text{partitions of }n\text{ into odd parts}\right\}
\rightarrow\left\{  \text{partitions of }n\text{ into distinct parts}\right\}.
\]
Given a partition $\lambda$ of $n$ into odd parts, we repeatedly merge pairs
of equal parts in $\lambda$ until no more equal parts appear. The final
result is $A\left(  \lambda\right)  $.

More precisely: if the original partition $\lambda$ contains an odd part $k$
exactly $m$ times, and if the binary representation of $m$ is
$m = \sum_{j=0}^{i} m_j 2^j$ with $m_j \in \{0,1\}$, then the partition
$A(\lambda)$ contains the number $2^j k$ exactly $m_j$ times for each $j$.
Since the binary digits $m_j$ are all $\leq 1$, the partition $A(\lambda)$ has
distinct parts.

The inverse map $B$ starts with a partition $\lambda$ into distinct parts.
Each part $k \cdot 2^i$ (with $k$ odd and $i \geq 0$) is replaced by $2^i$
copies of~$k$. Equivalently, one repeatedly splits even parts into halves
until only odd parts remain. The result $B(\lambda)$ is a partition into odd parts.

The maps $A$ and $B$ are mutually inverse, establishing the bijection.
Thus $p_{\operatorname{odd}}(n) = p_{\operatorname{dist}}(n)$.

In the formalization, this is proved using Glaisher's theorem from Mathlib,
which generalizes this result: for any positive integer $d$, the number of
partitions with parts not divisible by $d$ equals the number of partitions
where no part is repeated $d$ or more times. Euler's identity is the
special case $d = 2$.
\end{proof}

\subsection{Partitions with a given largest part}

\begin{lemma}
\label{lem.pars.transpose-def}
\lean{Nat.Partition.transpose}
\leanhelper
\leanok
For any partition $\lambda = (\lambda_1, \lambda_2, \ldots, \lambda_k)$,
the \emph{conjugate} (or \emph{transpose}) $\lambda^t$ is the partition
whose Young diagram is obtained by transposing the Young diagram of $\lambda$.
Explicitly, $\lambda^t = (\mu_1, \mu_2, \ldots, \mu_p)$ where $p$ is the
largest part of $\lambda$ and
\[
\mu_i = \left( \text{\# of parts of } \lambda \text{ that are } \geq i \right)
\quad \text{for each } i \in \{1, 2, \ldots, p\}.
\]
This satisfies $|\lambda^t| = |\lambda|$.
\end{lemma}

\begin{proof}
\leanok
The sum of the new parts equals the original size by a double-counting argument:
$\sum_i \mu_i = \sum_i \#\{j : \lambda_j \geq i\} = \sum_j \lambda_j = n$.
\end{proof}

\begin{lemma}
\label{lem.pars.transpose-size}
\lean{Nat.Partition.transpose_size}
\leanhelper
\leanok
The transpose preserves size: $|\lambda^t| = |\lambda|$,
i.e., if $\lambda$ is a partition of $n$, then $\lambda^t$ is also a partition of $n$.
\end{lemma}

\begin{proof}
\leanok
This is a restatement of the defining property of the transpose.
\end{proof}

\begin{lemma}
\label{lem.pars.transpose-involution}
\lean{Nat.Partition.transpose_transpose}
\leanhelper
\leanok
The transpose is an involution: $(\lambda^t)^t = \lambda$ for every partition $\lambda$.
\end{lemma}

\begin{proof}
\leanok
The proof uses the sorted list representation and the key combinatorial fact that
for a sorted decreasing list, the transpose operation is controlled by the
count of parts exceeding each threshold (Lemma~\ref{lem.pars.sorted-countP-gt-iff}).
This establishes that applying the Young diagram transpose twice recovers the original partition.
\end{proof}

\begin{lemma}
\label{lem.pars.sorted-countP-gt-iff}
\lean{Nat.Partition.sorted_countP_gt_iff}
\leanhelper
\leanok
For a sorted decreasing list $\ell = (\ell_0, \ell_1, \ldots)$,
indices $j < |\ell|$, and any $i\in\mathbb{N}$:
\[
|\{k : \ell_k > i\}| > j \iff \ell_j > i.
\]
This is the key combinatorial lemma underlying the Young diagram transpose involution:
it says the number of parts exceeding~$i$ is $>j$ precisely when
the $j$-th part itself exceeds~$i$.
\end{lemma}

\begin{proof}
\textbf{Forward:} If $|\{k : \ell_k > i\}| \le j$, then $\ell_j \le i$ (since
the sorted property implies all elements from position $j$ onward are $\le \ell_j$,
so at most $j$ elements exceed~$i$). Contrapositive gives the result.

\textbf{Backward:} If $\ell_j > i$, then all $\ell_0, \ldots, \ell_j$ are $> i$
(by the sorted property), so $|\{k : \ell_k > i\}| \ge j+1 > j$.
\end{proof}

\begin{lemma}
\label{lem.pars.transpose-length-eq-largestPart}
\lean{Nat.Partition.transpose_length_eq_largestPart}
\leanhelper
\leanok
The length of the transpose equals the largest part of the original partition:
$\ell(\lambda^t) = $ largest part of $\lambda$.
\end{lemma}

\begin{proof}
\leanok
The transpose has as many parts as there are columns in the Young diagram,
which equals the number of boxes in the first row, i.e., the largest part.
\end{proof}

\begin{lemma}
\label{lem.pars.filter-card-pos-of-lt-largest}
\lean{Nat.Partition.filter_card_pos_of_lt_largest}
\leanhelper
\leanok
For $i < $ largest part of $p$, the set of parts of $p$ that are $> i$ is nonempty.
\end{lemma}

\begin{proof}
If all parts were $\le i$, then the largest part would be $\le i$, contradicting
$i <$ largest part.
\end{proof}

\begin{lemma}
\label{lem.pars.filter-gt-zero-card-eq}
\lean{Nat.Partition.filter_gt_zero_card_eq}
\leanhelper
\leanok
For a partition $p$, filtering to keep only the positive parts preserves the cardinality
(since all parts are already positive).
\end{lemma}

\begin{proof}
Every part of a partition is positive, so restricting to positive parts does not change the count.
\end{proof}

\begin{lemma}
\label{lem.pars.transpose-largestPart-eq-length}
\lean{Nat.Partition.transpose_largestPart_eq_length}
\leanhelper
\leanok
The largest part of the transpose equals the length of the original partition:
largest part of $\lambda^t$ = $\ell(\lambda)$.
\end{lemma}

\begin{proof}
\leanok
The largest part of the transpose is the number of rows in the transposed diagram,
which equals the number of rows of the original, i.e., $\ell(\lambda)$.
\end{proof}

\begin{proposition}
\label{prop.pars.pkn=dual}
\lean{Nat.Partition.partsCount_eq_largestPartCount}
\leantarget
\leanok
Let $n\in\mathbb{N}$ and $k\in\mathbb{N}$. Then,%
\[
p_{k}\left(  n\right)  =\left(  \text{\# of partitions of }n\text{ whose
largest part is }k\right)  .
\]
\end{proposition}

\begin{proof}
We use the transpose as a bijection between:
\begin{itemize}
\item partitions of $n$ with $k$ parts, and
\item partitions of $n$ with largest part $k$.
\end{itemize}

For any partition $\lambda=\left(
\lambda_{1},\lambda_{2},\ldots,\lambda_{k}\right)  $, we define the
\emph{conjugate} $\lambda^{t}$ by $\lambda^{t}=\left(  \mu_{1}%
,\mu_{2},\ldots,\mu_{p}\right)  $, where $p$ is the largest part of $\lambda$
and $\mu_{i}=\left(  \text{\# of parts of }\lambda\text{ that are }\geq i\right)
$ for each $i\in\left\{  1,2,\ldots,p\right\}  $.
We have $\left\vert \lambda^{t}\right\vert =\left\vert \lambda\right\vert $ and
$\left(  \lambda^{t}\right)  ^{t}=\lambda$ for each partition $\lambda$, and
the largest part of $\lambda^{t}$ equals the length of $\lambda$.
By Lemma~\ref{lem.pars.transpose-largestPart-eq-length}, the transpose sends
a partition with $k$ parts to one whose largest part is $k$.
By Lemma~\ref{lem.pars.transpose-involution}, the transpose is an involution,
hence a bijection. This establishes the result.
\end{proof}

\begin{corollary}
\label{cor.pars.p0kn=dual}
\lean{Nat.Partition.partsCount_sum_eq_largestPartLeq}
\leantarget
\leanok
Let $n\in\mathbb{N}$ and $k\in\mathbb{N}$. Then,%
\begin{align*}
&  p_{0}\left(  n\right)  +p_{1}\left(  n\right)  +\cdots+p_{k}\left(
n\right) \\
&  =\left(  \text{\# of partitions of }n\text{ whose largest part is }\leq
k\right)  .
\end{align*}
\end{corollary}

\begin{proof}
We have%
\begin{align*}
&  p_{0}\left(  n\right)  +p_{1}\left(  n\right)  +\cdots+p_{k}\left(
n\right)
=\sum_{i=0}^{k}p_{i}\left(  n\right)
=\sum_{i=0}^{k}\left(  \text{\# of partitions of }n\text{ whose largest
part is }i\right)
\end{align*}
(by Proposition~\ref{prop.pars.pkn=dual} applied to each $i$). This sum counts
partitions of $n$ whose largest part is $\leq k$.
\end{proof}

\begin{lemma}
\label{lem.pars.partsCount-sum-eq-partsLeqCount}
\lean{Nat.Partition.partsCount_sum_eq_partsLeqCount}
\leanhelper
\leanok
The sum $p_0(n) + p_1(n) + \cdots + p_m(n)$ equals the count of partitions
with all parts $\leq m$:
\[
\sum_{k=0}^{m} p_k(n) = p_{\operatorname{parts}\leq m}(n).
\]
This combines Corollary~\ref{cor.pars.p0kn=dual} with the equivalence
``largest part $\leq m$'' $\iff$ ``all parts $\leq m$''
(Lemma~\ref{lem.pars.largestPart-le-iff}).
\end{lemma}

\begin{proof}
Direct from the two referenced results.
\end{proof}

\begin{theorem}
\label{thm.pars.main-gf-0n}
\lean{Nat.Partition.partsCountSum_genFun}
\leantarget
\leanok
Let $m\in\mathbb{N}$. Then,%
\[
\sum_{n\in\mathbb{N}}\left(  p_{0}\left(  n\right)  +p_{1}\left(  n\right)
+\cdots+p_{m}\left(  n\right)  \right)  x^{n}=\prod_{k=1}^{m}\dfrac{1}%
{1-x^{k}}.
\]
\end{theorem}

\begin{proof}
For each $n\in\mathbb{N}$, Corollary
\ref{cor.pars.p0kn=dual} shows that%
\[
p_{0}\left(  n\right)  +p_{1}\left(  n\right)  +\cdots+p_{m}\left(
n\right)
=\left(  \text{\# of partitions of }n\text{ whose largest part is }\leq
m\right)
=p_{\operatorname{parts}\leq m}\left(  n\right)  ,
\]
where $p_{\operatorname{parts}\leq m}\left(  n\right)  $ is defined as in
Theorem \ref{thm.pars.main-gf-parts-n}. Hence,%
\[
\sum_{n\in\mathbb{N}}\left(  p_{0}\left(  n\right)  +p_{1}\left(
n\right)  +\cdots+p_{m}\left(  n\right)  \right)  x^{n}=\sum_{n\in\mathbb{N}}p_{\operatorname{parts}%
\leq m}\left(  n\right)  x^{n}=\prod_{k=1}^{m}\dfrac{1}{1-x^{k}}
\]
(by Theorem \ref{thm.pars.main-gf-parts-n}).
\end{proof}

\subsection{Counting partitions by parts and largest part}

\begin{definition}
\label{def.pars.partsAndLargestCount}
\lean{Nat.Partition.partsAndLargestCount, Nat.Partition.partsAndLargestCountTotal}
\leanhelper
\leanok
Let $k,\ell\in\mathbb{N}$.
\begin{enumerate}
\item $q_{k,\ell}(n) := |\{p \vdash n : \ell(p) = k,\;\text{largest part of }p = \ell\}|$.
\item $Q(k,\ell) := \sum_{n=k}^{k\ell} q_{k,\ell}(n)$, the total number of such partitions over all valid sizes.
\end{enumerate}
\end{definition}

\begin{proposition}
\label{prop.pars.qbinom.intro-count-binom}
\lean{Nat.Partition.partsAndLargestCountTotal_eq}
\leantarget
\leanok
For positive integers $k$ and $\ell$,
\[
Q(k,\ell) = \binom{k+\ell-2}{k-1}.
\]
\end{proposition}

\begin{proof}
\leanok
The proof uses a bijection between the disjoint union
$\bigsqcup_{n\in [k, k\ell]} \{p \vdash n : \ell(p)=k,\;\text{largest}=\ell\}$
and the set of $(k-1)$-element multisubsets of $\{0,1,\ldots,\ell-1\}$, whose cardinality is
$\binom{\ell + (k-1) - 1}{k-1} = \binom{k+\ell-2}{k-1}$ by the stars-and-bars formula.

\textbf{Forward map:} A $(k-1)$-element multisubset $\{a_1,\ldots,a_{k-1}\}\subseteq\{0,1,\ldots,\ell-1\}$
maps to the partition $(\ell,\, a_1+1,\, \ldots,\, a_{k-1}+1)$ sorted in decreasing order.
This has $k$ parts, largest part $\ell$, and size $\ell + \sum(a_i+1) \in [k, k\ell]$.

\textbf{Backward map:} A partition $(\ell, \lambda_2,\ldots,\lambda_k)$ maps to the multiset
$\{\lambda_2-1, \ldots, \lambda_k-1\}\subseteq\{0,1,\ldots,\ell-1\}$,
since each $\lambda_i\in\{1,\ldots,\ell\}$.

The maps are mutual inverses:
incrementing each element of the multiset by~$1$ after decrementing undoes the decrement on parts $\ge 1$,
and prepending $\ell$ to the result of erasing $\ell$ recovers the original multiset.
\end{proof}

\begin{lemma}
\label{lem.pars.symToPartsMultiset-props}
\lean{Nat.Partition.symToPartsMultiset_pos, Nat.Partition.symToPartsMultiset_card, Nat.Partition.symToPartsMultiset_fold_max, Nat.Partition.symToPartsMultiset_erase, Nat.Partition.symToPartsMultiset_sum_range}
\leanhelper
\leanok
Properties of the multiset $(\ell, a_1+1, a_2+1, \ldots, a_{k-1}+1)$
for a $(k-1)$-element multisubset $(a_1, \ldots, a_{k-1})$ of $\{0,1,\ldots,\ell-1\}$:
\begin{enumerate}
\item All elements are positive (since $\ell\ge 1$ and each $a_i+1\ge 1$).
\item Cardinality is $k$ (one $\ell$ plus $k-1$ from the multisubset).
\item Largest part is $\ell$.
\item Removing $\ell$ recovers the remaining parts $(a_1+1, \ldots, a_{k-1}+1)$.
\item The sum lies in $[k, k\ell]$.
\end{enumerate}
\end{lemma}

\begin{proof}
\leanok
Each property follows from elementary multiset operations. In particular,
(3) uses $a_i+1\le\ell$ for $a_i \in \{0,1,\ldots,\ell-1\}$, and
(5) uses the bounds $1\le a_i+1\le\ell$.
\end{proof}

\subsection{Partition number vs.\ sums of divisors}

\begin{lemma}
\label{lem.pars.lhs-eq-sum-parts}
\lean{Nat.Partition.lhs_eq_sum_parts}
\leanhelper
\leanok
For any set $I\subseteq\{1,2,3,\ldots\}$ and $n\in\mathbb{N}$:
\[
n\cdot |\operatorname{restricted}(n, I)| = \sum_{p\in\operatorname{restricted}(n,I)} |\,p\,|.
\]
\end{lemma}

\begin{proof}
\leanok
Each partition $p$ of $n$ has size $|p| = n$, so the sum of a constant
equals the constant times the count.
\end{proof}

\begin{lemma}
\label{lem.pars.sum-count-eq-sum-card-ge}
\lean{Nat.Partition.sum_count_eq_sum_card_ge}
\leanhelper
\leanok
For any finite set $s$ of partitions of $n$ and any $d\in\mathbb{N}$:
\[
\sum_{p\in s} \operatorname{count}(d, p) = \sum_{j=1}^{n} |\{p\in s : \operatorname{count}(d,p)\ge j\}|.
\]
\end{lemma}

\begin{proof}
\leanok
This is the standard identity $\sum_i f(i) = \sum_{j\ge 1}|\{i : f(i)\ge j\}|$
applied with $f(p) = \operatorname{count}(d, p)$, bounded by~$n$ since at most $n/d$ copies
of~$d$ fit in a partition of~$n$.
\end{proof}

\begin{lemma}
\label{lem.pars.sum-parts-as-weighted-count}
\lean{Nat.Partition.sum_parts_as_weighted_count}
\leanhelper
\leanok
For any partition $p$ of $n$,
\[
\text{(sum of parts of $p$)} = \sum_{d \text{ distinct in } p} d \cdot \operatorname{count}(d, p).
\]
\end{lemma}

\begin{proof}
\leanok
By induction on the multiset of parts.
\end{proof}

\begin{lemma}
\label{lem.pars.count-le-div}
\lean{Nat.Partition.count_le_div_of_partition}
\leanhelper
\leanok
For any partition $p$ of $n$ and $d > 0$,
$\operatorname{count}(d, p) \leq n/d$.
\end{lemma}

\begin{proof}
\leanok
At most $n/d$ copies of $d$ can fit in a partition of $n$.
\end{proof}

\begin{lemma}
\label{lem.pars.removePartCopies}
\lean{Nat.Partition.removePartCopies, Nat.Partition.addPartCopies}
\leanhelper
\leanok
If a partition $p$ of $n$ contains $j$ or more copies of a part $d$,
then removing $j$ copies of $d$ yields a valid partition of $n - d \cdot j$.
Conversely, adding $j$ copies of a positive integer $d$ to a partition of $m$
yields a partition of $m + d \cdot j$.
\end{lemma}

\begin{proof}
\leanok
Removing copies preserves positivity of remaining parts and adjusts the sum.
Adding copies preserves positivity since $d > 0$.
\end{proof}

\begin{lemma}
\label{lem.pars.card-count-ge-eq-restricted}
\lean{Nat.Partition.card_count_ge_eq_restricted}
\leanhelper
\leanok
Let $I \subseteq \{1, 2, 3, \ldots\}$, let $d \in I$, and let $d \cdot j \leq n$.
Then
\[
\#\{p \in \operatorname{restricted}(n, I) : \operatorname{count}(d, p) \geq j\}
= \#\operatorname{restricted}(n - d \cdot j, I).
\]
The bijection sends a partition $p$ to the partition obtained by removing $j$
copies of $d$, and its inverse adds $j$ copies of $d$.
\end{lemma}

\begin{proof}
The forward map removes $j$ copies of $d$ from the partition.
The backward map adds $j$ copies. Both preserve the property that all parts are in $I$.
Injectivity and surjectivity are verified by showing the maps are mutual inverses.
\end{proof}

\begin{lemma}
\label{lem.pars.divisor-sum-reindex}
\lean{Nat.Partition.divisor_sum_reindex}
\leanhelper
\leanok
For any set $I$ of positive integers, any $n \in \mathbb{N}$, and any function $f$:
\[
\sum_{k=0}^{n-1} \sum_{\substack{d \mid k+1 \\ d \in I}} d \cdot f(n - k - 1)
= \sum_{\substack{d \in I \\ d \leq n}} d \cdot \sum_{j=1}^{\lfloor n/d \rfloor} f(n - d \cdot j).
\]
The bijection is $(k, d) \leftrightarrow (d, j)$ where $k = d \cdot j - 1$.
\end{lemma}

\begin{proof}
The reindexing bijection is $(k, d) \leftrightarrow (d, j)$ where $k = d \cdot j - 1$.
The forward map sends $(k, d)$ to $(d, (k+1)/d)$ and the backward map sends $(d, j)$ to $(d \cdot j - 1, d)$.
\end{proof}

\begin{theorem}
\label{thm.pars.sigma1-I}
\lean{Nat.Partition.partitionCount_divisorSum_restricted}
\leantarget
\leanok
Let $I$ be a subset of $\left\{  1,2,3,\ldots
\right\}  $. For each $n\in\mathbb{N}$, let $p_{I}\left(  n\right)  $ be the
\# of partitions $\lambda$ of $n$ such that all parts of $\lambda$ belong to
$I$.

For any positive integer $n$, let $\sigma_{I}\left(  n\right)  $ denote the
sum of all positive divisors of $n$ that belong to $I$.

For any $n\in\mathbb{N}$, we have%
\[
np_{I}\left(  n\right)  =\sum_{k=1}^{n}\sigma_{I}\left(  k\right)
p_{I}\left(  n-k\right)  .
\]
\end{theorem}

\begin{proof}
The proof uses a double counting argument. Both sides count weighted pairs.

\textbf{LHS Analysis:}
\[
n \cdot |{\operatorname{restricted}(n, I)}|
= \sum_{p} |p|
= \sum_{p} \sum_{\substack{d \text{ distinct}\\\text{in } p}} d \cdot \operatorname{count}(d, p)
\]
Swapping the order of summation and using the identity
$\sum_p \operatorname{count}(d, p) = \sum_{j \geq 1} |\{p : \operatorname{count}(d, p) \geq j\}|$,
together with the bijection from Lemma~\ref{lem.pars.card-count-ge-eq-restricted}, gives:
\[
= \sum_{\substack{d \in I \\ d \leq n}} d \cdot \sum_{j=1}^{\lfloor n/d \rfloor}
|\operatorname{restricted}(n - d \cdot j, I)|.
\]

\textbf{RHS Analysis:}
\[
\sum_{k=0}^{n-1} \sigma_I(k+1) \cdot |\operatorname{restricted}(n-k-1, I)|
= \sum_{k=0}^{n-1} \sum_{\substack{d \mid k+1 \\ d \in I}} d \cdot |\operatorname{restricted}(n-k-1, I)|.
\]
By the reindexing lemma (Lemma~\ref{lem.pars.divisor-sum-reindex}), this equals
the same expression.
\end{proof}

\begin{theorem}
\label{thm.pars.sigma1}
\lean{Nat.Partition.partitionCount_divisorSum}
\leantarget
\leanok
For any positive integer $n$, let $\sigma\left(
n\right)  $ denote the sum of all positive divisors of $n$. (For example,
$\sigma\left(  6\right)  =1+2+3+6=12$ and $\sigma\left(  7\right)  =1+7=8$.)

For any $n\in\mathbb{N}$, we have%
\[
np\left(  n\right)  =\sum_{k=1}^{n}\sigma\left(  k\right)  p\left(
n-k\right)  .
\]
\end{theorem}

\begin{proof}
This is the particular case of Theorem~\ref{thm.pars.sigma1-I} for $I=\left\{  1,2,3,\ldots\right\}$.
We specialize the restricted version with $I = \{n \in \mathbb{N} : n > 0\}$.
Since all parts of a partition are positive, the restricted set equals the set of all partitions,
so $p_I(n) = p(n)$. Similarly, $\sigma_I(n) = \sigma(n)$ since all divisors of a positive integer
are positive.
\end{proof}

\subsubsection{Generating function proof of the recurrence (alternative approach)}

The divisor sum recurrence can also be proved via generating functions.
Define $P := \sum_{n\ge 0} p(n) x^n$ and $S := \sum_{k\ge 1} \sigma(k) x^k$.
Then $X\cdot P' = S\cdot P$, and comparing coefficients gives the result.

\begin{definition}
\label{def.pars.genFun-P-S}
\lean{Nat.Partition.P, Nat.Partition.S}
\leanhelper
\leanok
The partition generating function $P := \sum_{n\ge 0} p(n) x^n \in \mathbb{Z}[[x]]$
and the divisor sum generating function $S := \sum_{k\ge 1} \sigma(k) x^k \in \mathbb{Z}[[x]]$.
\end{definition}

\begin{lemma}
\label{lem.pars.coeff-P}
\lean{Nat.Partition.coeff_P}
\leanhelper
\leanok
The coefficient of $x^n$ in $P$ equals $p(n)$: $[x^n]P = p(n)$.
\end{lemma}

\begin{proof}
\leanok
Immediate from the definition of $P$.
\end{proof}

\begin{lemma}
\label{lem.pars.coeff-X-mul-derivative-P}
\lean{Nat.Partition.coeff_X_mul_derivative_P}
\leanhelper
\leanok
The coefficient of $x^n$ in $X\cdot P'$ equals $n\cdot p(n)$:
$[x^n](X \cdot P') = n \cdot [x^n]P$.
\end{lemma}

\begin{proof}
\leanok
Since $[x^n](X\cdot f') = n\cdot[x^n]f$ for any FPS~$f$.
\end{proof}

\begin{lemma}
\label{lem.pars.coeff-S-mul-P}
\lean{Nat.Partition.coeff_S_mul_P}
\leanhelper
\leanok
The coefficient of $x^n$ in $S\cdot P$ equals the RHS of the recurrence:
\[
[x^n](S\cdot P) = \sum_{k=1}^{n} \sigma(k)\cdot p(n-k).
\]
\end{lemma}

\begin{proof}
\leanok
Expand the Cauchy product $[x^n](S\cdot P) = \sum_{i+j=n} [x^i]S\cdot [x^j]P$,
noting that $[x^0]S = 0$ so the $i=0$ term vanishes.
\end{proof}

\begin{theorem}
\label{thm.pars.X-mul-derivative-P-eq-S-mul-P}
\lean{Nat.Partition.X_mul_derivative_P_eq_S_mul_P}
\leanhelper
\leanok
The generating function identity $X\cdot P' = S\cdot P$ holds as an equality of
formal power series in $\mathbb{Z}[[x]]$.
\end{theorem}

\begin{proof}
\leanok
Comparing coefficients at each $x^n$:
$[x^n](X\cdot P') = n\cdot p(n) = \sum_{k=1}^{n}\sigma(k)\cdot p(n-k) = [x^n](S\cdot P)$,
where the middle equality is Theorem~\ref{thm.pars.sigma1}.
\end{proof}
